\clearpage
\section{Заключение}

Для области варианта 14 решена задача Дирихле для уравнения Пуассона методом фиктивных областей.
Коэффициенты разностной схемы учитывают долю площади каждой ячейки внутри области.
Система решается методом сопряжённых градиентов с диагональным предобуславливанием; замена неизвестной на поправку к $\sin(x+y)$ устраняет вычитание больших штрафных слагаемых.

Последовательная, OpenMP-, MPI- и гибридная реализации дают согласованные результаты.
Проверка с известным гладким решением показывает уменьшение ошибки примерно в четыре раза при двукратном сгущении сетки.
Для исходной задачи получено поле решения на сетке $800\times1200$.

На IBM Polus выполнены все 37 конфигураций основного набора и приложения, включая последовательные базовые расчёты.
На сетке $800\times1200$ MPI с 32 процессами на двух узлах сокращает медианное время с $28.095$ до $1.090$ с, что соответствует ускорению $25.779$.
Гибридная версия на 32 физических ядрах даёт ускорение $15.933$; OpenMP с 32 аппаратными потоками на 16 ядрах одного узла даёт $12.779$.

Для части MPI-конфигураций увеличение числа процессов приводит к повторяемому замедлению.
Причины замедлений в проведённых измерениях не установлены.
Полученные времена и ускорения относятся к зафиксированным условиям данной серии запусков.
